% Conclusión individual: Mera
Mi parte de la práctica fue la implementación con gRPC: escribí el contrato \texttt{messaging.proto}, el servidor y el cliente, y los integré con el reloj de Lamport de Roger. Una decisión bastanete importante fue que mi cliente tuviera la misma interfaz que el de TCP (\texttt{conectar}, \texttt{enviar}, \texttt{cerrar}), para que el compañero Mario pudiera medir los dos mecanismos con el mismo script. También decidí abrir un solo canal y esperar a que estuviera listo antes de medir, para no contar la conexión dentro de la latencia.

El problema que más me costó fue que el servidor seguía respondiendo con 6 en vez de 7 después de integrar el reloj; el código estaba bien, pero el archivo no se había guardado. En los 20 intercambios obtuve la misma secuencia de relojes que en TCP (acuse final 79 y reloj local 80), lo que me confirmó que el transporte no cambia el orden de Lamport. De esa manera aprendí que el tipado de gRPC facilita el trabajo, pero no reemplaza la validación: una solicitud sin \texttt{timestamp} llega con valor 0 sin dar error. Con 100 bytes, gRPC fue unas 4,5 veces más lento que TCP, así que en el PFC lo usaría donde importe más un contrato claro entre servicios que la latencia mínima.